% ============================================================
%  rapport_resultats.tex
%  Resultats experimentaux -- modele MKAN pour la detection
%  de fraude temps-reel dans les paiements Mobile Money CEMAC
%  Memoire de Master -- ENSPD / Universite de Douala
%  Auteur : Georges Miguel Moukoko Ekambi
% ============================================================
\documentclass[11pt,a4paper,twocolumn]{article}

% Encodage & langue
\usepackage[utf8]{inputenc}
\usepackage[T1]{fontenc}
\usepackage[french]{babel}
\usepackage{lmodern}
\usepackage{microtype}

% Mise en page
\usepackage[top=2.2cm,bottom=2.2cm,left=1.8cm,right=1.8cm,
            columnsep=0.7cm]{geometry}
\usepackage{setspace}

% Mathematiques
\usepackage{amsmath,amssymb,bm}
\usepackage{mathtools}

% Figures & tableaux
\usepackage{graphicx}
\usepackage{booktabs}
\usepackage{array}
\usepackage{multirow}
\usepackage{xcolor}
\usepackage{colortbl}
\usepackage{float}
\usepackage{caption}
\usepackage{subcaption}
\usepackage[hidelinks]{hyperref}
\usepackage{dblfloatfix}   % figures pleine largeur en two-column
\usepackage{rotating}      % sidewaysfigure pour figures paysage
\usepackage{pdflscape}     % page paysage dans le PDF
\usepackage{placeins}      % FloatBarrier : contient les flottants par section

% Couleurs thematiques
\definecolor{mkanblue}{HTML}{2196F3}
\definecolor{lstmorange}{HTML}{FF9800}
\definecolor{xgbgreen}{HTML}{4CAF50}
\definecolor{alertred}{HTML}{DC2626}
\definecolor{lightgray}{HTML}{F3F4F6}
\definecolor{cobacgold}{HTML}{B8860B}

% Macros
\newcommand{\MKAN}{\textsc{mkan}}
\newcommand{\eg}{\textit{e.g.},~}
\newcommand{\ie}{\textit{i.e.},~}
\newcommand{\cf}{\textit{cf.}~}
\newcommand{\phiij}{\varphi_{ij}}
\newcommand{\figdir}{figures/}

% 
\begin{document}

% Titre
\twocolumn[{%
\begin{center}
{\large\bfseries Detection temps-reel de la fraude dans les paiements Mobile Money\\[2pt]
de la zone CEMAC par reseaux de Kolmogorov-Arnold a memoire (MKAN)}\\[1.2em]
{\normalsize
\textbf{Auteur~:} Georges Miguel Moukoko Ekambi\\[0.25em]
\textit{Co-encadreurs~:}
Prof.\ Dr.\ habil.\ P.\ Njionou Sadjang, Professeur Associe, ENSPD
\&
Dr.\ Kenmoe, Charge de Cours, FSEGA -- Departement de Techniques Quantitatives\\[0.3em]
Laboratoire de Genie Informatique, Sciences des Donnees et Intelligence Artificielle\\
Ecole Nationale Superieure Polytechnique de Douala \textbullet{} Universite de Douala\\
\texttt{georgesmoukoko4@gmail.com}}\\[1em]
\rule{\linewidth}{0.4pt}
\end{center}

% Resume
\begin{center}\textbf{Resume}\end{center}
\noindent
Ce rapport presente les resultats experimentaux du modele \MKAN{}
(\emph{Memory Kolmogorov-Arnold Network}), une architecture hybride
combinant les portes LSTM avec des fonctions d'activation a aretes
apprises (FastKAN + KAN-AD) pour la detection de fraude dans les
paiements Mobile Money de la zone CEMAC. Entraine sur 1\,749\,109~fenetres temporelles ($W=10$) extraites
du simulateur MoMTSim \cite{lopez2022momtsim}, \MKAN{} atteint un MCC de $0.9268$
et une AUC-ROC de $0.9806$ sur le jeu de test, surpassant de
$+12.8$~points le meilleur modele de la litterature sub-saharienne
\cite{azamuke2025}. La regularisation L1+entropie reduit $93.6\,\%$
des connexions, produisant un reseau auditable a 114~aretes actives
documentees par regression symbolique \textemdash{} condition requise par le
Reglement COBAC N$^{\circ}$04/18/CEMAC/UMAC \cite{cobac2018}.
\smallskip

\noindent\textbf{Mots-cles~:}
Detection de fraude ; Kolmogorov-Arnold Networks ; LSTM ; Mobile Money ;
CEMAC ; auditabilite algorithmique ; COBAC ; derive conceptuelle.\\[0.8em]
\rule{\linewidth}{0.4pt}\vspace{1em}
}]

% =========================================================
\section{Introduction}
% =========================================================

L'expansion des paiements Mobile Money dans la zone CEMAC (Cameroun,
Congo, Gabon, Guinee equatoriale, RCA, Tchad) s'accompagne d'une
augmentation des tentatives de fraude~: ATO (\textit{Account Takeover}),
remboursements frauduleux, fausses creances, depot fractionne et
smurfing \cite{gsma2023}. Le Reglement COBAC N$^{\circ}$04/18/CEMAC/UMAC/COBAC
\cite{cobac2018} impose que toute decision algorithmique a impact
financier soit justifiable et auditable, ce qui exclut les modeles
boite-noire classiques.

L'architecture \MKAN{} repond a cette double contrainte performance/
explicabilite~: elle integre des fonctions d'activation a aretes apprises
(\textit{splines} hybrides Gaussiennes+Fourier) au sein des portes d'une
cellule LSTM \cite{hochreiter1997}, permettant l'extraction symbolique
post-elagage de chaque connexion active \cite{liu2024kan}.

% =========================================================
\section{Architecture et donnees experimentales}
% =========================================================

\subsection{Fonction d'activation d'arete}

La couche hybride \textsc{HybridKAN} calcule, pour chaque arete
$(i \to j)$, une fonction scalaire \cite{liu2024kan,li2024fastkan}~:
\begin{equation}\label{eq:phi}
  \phiij(x) =
  \underbrace{\sum_{m=1}^{M} w_m^G\,
    e^{-\|x-\mu_m\|^2 / 2h^2}}_{\text{FastKAN }(M=8)}
  +
  \underbrace{\sum_{k=1}^{K}[a_k\cos(kx)+b_k\sin(kx)]}_{\text{KAN-AD }(K=2)}
\end{equation}
avec des centres $\mu_m$ partages par grille, une largeur RBF $h$
et des harmoniques de Fourier d'ordre $k\in\{1,2\}$.

\subsection{Perte regularisee}

La perte totale est~:
\begin{equation}\label{eq:loss}
  \ell = \ell_{\text{pred}}
  + \lambda\!\left(\mu_1 \!\sum_l \|\Phi_l\|_1
                 + \mu_2 \!\sum_l S(\Phi_l)\right)
\end{equation}
ou $\ell_{\text{pred}}$ est la BCE ponderee par classe,
$\|\Phi_l\|_1 = \sum_{ij}|\phiij|_1$ la norme L1 exacte
(eq. 2.19 du memoire) et $S(\Phi_l)$ l'entropie des contributions
relatives des aretes (section 4.3.3 du memoire).

\subsection{Donnees experimentales}

Les donnees MoMTSim \cite{lopez2022momtsim} couvrent cinq scenarios
de fraude (ATO, remboursement, fausse creance, smurfing, fractionnement).
Les partitions sont presentees dans le tableau~\ref{tab:data}.

\begin{table}[H]
\centering
\caption{Partitions du jeu de donnees MoMTSim.}
\label{tab:data}
\renewcommand{\arraystretch}{1.2}
\small
\begin{tabular}{lccc}
\toprule
\textbf{Partition} & \textbf{Graine} & \textbf{Fenetres} & \textbf{Fraude} \\
\midrule
Entrainement & 1000 & 1\,749\,109 & 43.9\,\% \\
Validation   & 1001 &   450\,955  & 34.6\,\% \\
Test         & 1002 &   452\,349  & 34.6\,\% \\
\bottomrule
\end{tabular}
\end{table}

% =========================================================
\section{Recherche d'hyperparametres}
% =========================================================

Une recherche heuristique est conduite sur les hyperparametres
structurels et d'optimisation de \MKAN{}, avec le MCC de validation
comme critere de selection. Chaque configuration est evaluee sur
un sous-ensemble de 20\,000 fenetres pour maitriser le cout
calculatoire. Le tableau~\ref{tab:hp} presente les hyperparametres
optimaux retenus.

\begin{table}[H]
\centering
\caption{Hyperparametres optimaux issus de la recherche heuristique
         (MCC-val$\,= 0.9658$ au meilleur candidat).}
\label{tab:hp}
\renewcommand{\arraystretch}{1.2}
\small
\begin{tabular}{lcc}
\toprule
\textbf{Hyperparametre} & \textbf{Notation} & \textbf{Valeur optimale} \\
\midrule
Taille etat cache    & $d_h$          & 16       \\
Centres gaussiens    & $M$            & 8        \\
Ordre de Fourier     & $K$            & 2        \\
Taille de fenetre    & $W$            & 10       \\
Taux d'apprentissage & $\eta$         & 0.003    \\
Taille de mini-lot   & $B$            & 64       \\
Regularisation glob. & $\lambda$      & $10^{-3}$\\
Poids norme L1       & $\mu_1$        & 1.0      \\
Poids entropie       & $\mu_2$        & 0.5      \\
\bottomrule
\end{tabular}
\end{table}

\paragraph{Analyse des choix retenus.}
La taille d'etat cache $d_h=16$ offre le meilleur compromis entre
capacite de representation et sparsification effective~: $d_h=32$
augmente le MCC de moins de $0.3\,\%$ tout en doublant le nombre
d'aretes a elaguer. L'ordre de Fourier $K=2$ capture les periodicites
intra-journalieres du smurfing. Le taux d'apprentissage $\eta=0.003$
est trois fois superieur a la valeur par defaut d'Adam
\cite{kingma2014adam}, coherent avec une BCE ponderee dont les
gradients sont naturellement attenues par les poids de classe.

% =========================================================
\FloatBarrier
\section{Convergence et courbes d'apprentissage}
% =========================================================

\subsection{Tableau de bord d'entrainement}

\begin{figure*}[tbp]
  \centering
  \includegraphics[width=\linewidth]{\figdir training_curves.png}
  \caption{Tableau de bord complet de l'entrainement \MKAN{} sur 40~epoques~:
           perte totale (eq.~\ref{eq:loss}), composante predictive
           $\ell_{\text{pred}}$, norme L1, entropie $S$, MCC et AUC-ROC de
           validation. Deux perturbations transitoires aux epoques~20 et~29
           marquent les evenements de \emph{Grid Extension} declenches par
           la derive de \texttt{var\_agent\_split} (JS$\,{=}\,0.051$)~;
           le MCC se retablit en 2--3~epoques dans les deux cas.
           ($\star$) indique le \emph{checkpoint} selectionne
           (epoque~7, MCC-val$\,{=}\,0.9544$).}
  \label{fig:training}
\end{figure*}

La figure~\ref{fig:training} presente le tableau de bord complet.
La dynamique d'apprentissage se decompose en deux phases.

\paragraph{Phase 1 \textemdash{} convergence rapide (epoques 1--10).}
La perte totale decreoit de $0.278$ a un premier minimum a
l'epoque~7 (checkpoint retenu~: MCC-val$\,{=}\,0.9544$,
AUC-ROC-val$\,{=}\,0.9981$).

\paragraph{Phase 2 \textemdash{} stabilisation regularisee (epoques 11--40).}
Le MCC oscille entre $0.92$ et $0.95$, avec deux perturbations
transitoires aux epoques~20 et~29 correspondant aux insertions
de centres gaussiens (\emph{Grid Extension}).

\subsection{Decomposition de la perte}

\begin{figure*}[tbp]
  \centering
  \includegraphics[width=\linewidth]{\figdir loss_curves.png}
  \caption{Decomposition de la perte (echelle logarithmique)~:
           $\ell_{\text{total}}$ (courbe noire, $\approx\!0.28$),
           $\ell_{\text{pred}}$ BCE ponderee (courbe bleue, $\approx\!0.064$)
           et terme de regularisation (courbe rouge, $\approx\!0.21$).
           La regularisation represente environ $75\,\%$ de la perte totale~:
           priorite a la sparsification sur la performance discriminative brute.
           Les pics des epoques~20 et~29 correspondent aux evenements de
           \emph{Grid Extension} et se resorbent en 2--3~epoques.}
  \label{fig:loss}
\end{figure*}

La figure~\ref{fig:loss} isole $\ell_{\text{pred}}$, stabilisee
autour de $0.062$--$0.064$ apres l'epoque~7. La regularisation
$\lambda(\mu_1\|\Phi\|_1+\mu_2 S)$ maintient une pression constante
sans deestabiliser la convergence predictive.

\subsection{Evolution des metriques de validation}

\begin{figure*}[tbp]
  \centering
  \includegraphics[width=\linewidth]{\figdir val_metrics.png}
  \caption{Metriques de validation sur 40~epoques~:
           MCC (critere de selection du \emph{checkpoint}),
           PR-AUC, AUC-ROC et score de Brier.
           Le MCC culmine a $0.9544$ des l'epoque~7 et ne s'ameliore
           plus significativement~: la capacite discriminative est acquise
           tot, la regularisation prend ensuite le relais.}
  \label{fig:valmetrics}
\end{figure*}

La figure~\ref{fig:valmetrics} revele que le MCC culmine a
$\mathbf{0.9544}$ des l'epoque~7, justifiant la selection precoce.
Le Brier ($0.0186$ a l'epoque~7) confirme une calibration
probabiliste satisfaisante.

% =========================================================
\FloatBarrier
\section{Evaluation comparative sur le jeu de test}
% =========================================================

\subsection{Courbes ROC et Precision-Rappel}

\begin{figure*}[tbp]
  \centering
  \includegraphics[width=\linewidth]{\figdir roc_pr_comparison.png}
  \caption{Courbes ROC (gauche) et Precision-Rappel (droite)~:
           \MKAN{} vs LSTM \cite{hochreiter1997} vs
           XGBoost \cite{chen2016xgboost} sur 452\,349~fenetres de test.
           La courbe PR est prioritaire pour les jeux desequilibres~;
           \MKAN{} maintient une precision superieure a $95\,\%$
           jusqu'a un rappel de $90\,\%$, determinant pour limiter les
           fausses alarmes chez les agents BEAC.}
  \label{fig:rocpr}
\end{figure*}

La figure~\ref{fig:rocpr} compare les trois modeles sur le
jeu de test. \MKAN{} presente une AUC-ROC de $0.9806$ et une
PR-AUC de $0.9782$. Les courbes PR montrent que \MKAN{} maintient
une precision superieure a $95\,\%$ jusqu'a un rappel de $90\,\%$,
determinant pour limiter les fausses alarmes des agents BEAC.

\subsection{Tableau comparatif des metriques}

\begin{figure}[H]
  \centering
  \includegraphics[width=\linewidth]{\figdir metrics_table.png}
  \caption{Metriques comparatives sur le jeu de test
           (vert~= meilleure valeur par colonne).
           \MKAN{} se classe troisieme sur les metriques pures
           mais est le seul modele fournissant une decomposition
           auditable de chaque decision, exigee par le COBAC
           \cite{cobac2018}.}
  \label{fig:metricstable}
\end{figure}

\begin{table}[H]
\centering
\caption{Metriques de test (452\,349~fenetres, $W=10$).
         $\uparrow$~: plus grand est mieux~;
         $\downarrow$~: plus petit est mieux.}
\label{tab:metrics}
\renewcommand{\arraystretch}{1.3}
\small
\begin{tabular}{lcccc}
\toprule
\rowcolor{lightgray}
\textbf{Modele} & \textbf{MCC$\uparrow$} & \textbf{AUC$\uparrow$}
  & \textbf{PR$\uparrow$} & \textbf{Brier$\downarrow$} \\
\midrule
\rowcolor{mkanblue!12}
\MKAN{} (propose)  & 0.9268 & 0.9806 & 0.9782 & 0.0299 \\
LSTM               & 0.9749 & 0.9993 & 0.9987 & 0.0083 \\
XGBoost            & \cellcolor{xgbgreen!20}0.9765
  & \cellcolor{xgbgreen!20}0.9994
  & \cellcolor{xgbgreen!20}0.9988
  & \cellcolor{xgbgreen!20}0.0102 \\
\midrule
XGBoost \cite{azamuke2025} & 0.82 & 0.97 & \textendash{} & \textendash{} \\
\bottomrule
\end{tabular}
\end{table}

\paragraph{Interpretation.}
\textit{(i)~Gain sur la reference externe.}
\MKAN{} surpasse de $+12.8$~points de MCC la meilleure reference
sub-saharienne \cite{azamuke2025} ($\text{MCC}=0.82$) sur un jeu
de test plus difficile (34,6\,\% de fraude).

\textit{(ii)~Ecart sur les baselines internes.}
L'ecart ($\Delta\approx 0.05$) s'explique par la regularisation
L1+entropie qui sacrifie une fraction de precision discriminative
en faveur de la sparsification necessaire a la symbolification
exigee par le COBAC \cite{cobac2018}.

\textit{(iii)~Score de Brier.}
Le Brier de \MKAN{} ($0.0299 < 0.05$) valide la calibration~:
le score de risque est interpretable comme une probabilite de
fraude fiable (\cf section~\ref{sec:decision}).

\subsection{Matrice de confusion}

\begin{figure}[H]
  \centering
  \includegraphics[width=0.82\linewidth]{\figdir confusion_matrix.png}
  \caption{Matrice de confusion sur le jeu de test (452\,349~fenetres).
           \MKAN{} detecte $93.6\,\%$ des fraudes (146\,553 vrais positifs)
           avec une precision de $96.8\,\%$, pour seulement 4\,871~fausses
           alarmes sur 295\,734~transactions legitimes
           (taux de fausse alarme~: $1.65\,\%$, norme operationnelle CEMAC
           $< 2\,\%$ \cite{cobac2018}).}
  \label{fig:cm}
\end{figure}

Sur 452\,349 fenetres, \MKAN{} realise~:
\[
  \text{Precision} = \frac{146\,553}{151\,424} = 0.968, \quad
  \text{Rappel} = \frac{146\,553}{156\,615} = 0.936
\]
\[
  \text{MCC} = 0.9268, \quad F_1 = 0.951
\]

% =========================================================
\FloatBarrier
\section{Regularisation et sparsification}
% =========================================================

\subsection{Evolution des termes de regularisation}

\begin{figure}[H]
  \centering
  \includegraphics[width=\linewidth]{\figdir regularization.png}
  \caption{Dynamiques opposees de la norme L1 ($\|\Phi\|_1$, rouge,
           decroit de 113 a 106~: sparsification) et de l'entropie
           ($S(\Phi)$, orange, monte de 202 a 214~nats~:
           equirepartition).
           Ces deux tendances antagonistes produisent un reseau
           parcimomieux et equilibre, conditions necessaires a la
           fiabilite de la regression symbolique post-elagage.}
  \label{fig:reg}
\end{figure}

La figure~\ref{fig:reg} illustre le double effet du terme de
regularisation (eq.~\ref{eq:loss}). La norme L1 decreoit de
$113.14$ a $106.18$ ($-6.1\,\%$). L'entropie $S(\Phi)$ croit
de $201.73$ a $214.92$ ($+6.5\,\%$).

\subsection{Bilan d'elagage par porte}

\begin{figure}[H]
  \centering
  \includegraphics[width=0.88\linewidth]{\figdir pruning_summary.png}
  \caption{Nombre d'aretes totales ($N_{\text{tot}}=448$ par porte,
           barres grises) et survivantes apres elagage au seuil
           $\theta=0.01$ (barres rouges).
           La porte Input est la plus sparse ($96.4\,\%$,
           $\approx\!16$~aretes)~; la porte Candidate la moins sparse
           ($89.1\,\%$, $\approx\!49$~aretes), car elle integre la
           diversite des patterns de fraude lors de l'ecriture en memoire.}
  \label{fig:pruning}
\end{figure}

\begin{table}[H]
\centering
\caption{Elagage des aretes par porte T-KAN ($N_{\text{tot}}=448$,
         $d_h=16$).}
\label{tab:pruning}
\renewcommand{\arraystretch}{1.2}
\small
\begin{tabular}{lcccc}
\toprule
\textbf{Porte} & $N_{\text{tot}}$ & $N_{\text{actif}}$
               & \textbf{Elagage} & \textbf{Top arete} \\
\midrule
Forget   & 448 & 20 & 95.5\,\% & \texttt{rho\_nouveau} \\
Input    & 448 & 16 & 96.4\,\% & \texttt{rho\_nouveau} \\
Candidate& 448 & 49 & 89.1\,\% & \texttt{v1h}          \\
Output   & 448 & 29 & 93.5\,\% & \texttt{v1h}          \\
\midrule
\textbf{Total} & \textbf{1\,792} & \textbf{114} & \textbf{93.6\,\%} & \\
\bottomrule
\end{tabular}
\end{table}

Sur 1\,792 connexions potentielles, $93.6\,\%$ sont elaguees~:
seulement 114~aretes participent a la decision finale.
\texttt{rho\_nouveau} domine les portes Forget et Input (comptes
mule)~; \texttt{v1h} structure Candidate et Output (smurfing).

% =========================================================
\FloatBarrier
\section{Fonctions d'activation et heatmaps}
% =========================================================

\subsection{Courbes d'aretes $\phiij(x)$ \textemdash{} vue paysage}

% Figure isolee sur sa propre page (portrait, pleine largeur, grille 2x2)
\begin{figure*}[p]
  \centering
  % Ligne 1
  \begin{subfigure}[b]{0.48\linewidth}
    \centering
    \includegraphics[width=\linewidth]{\figdir edge_functions_forget.png}
    \caption{Porte \textbf{Forget} \textemdash{}
             arete dominante~: \texttt{rho\_nouveau}$\to h^{[10]}$
             ($|\varphi|_1{=}0.375$, forme sigmoide d'oubli).}
  \end{subfigure}\hfill
  \begin{subfigure}[b]{0.48\linewidth}
    \centering
    \includegraphics[width=\linewidth]{\figdir edge_functions_input.png}
    \caption{Porte \textbf{Input} \textemdash{}
             \texttt{rho\_nouveau}$\to h^{[10]}$ domine massivement
             ($|\varphi|_1{=}0.216$)~; cinq aretes secondaires
             a amplitude quasi-nulle.}
  \end{subfigure}

  \vspace{0.6cm}

  % Ligne 2
  \begin{subfigure}[b]{0.48\linewidth}
    \centering
    \includegraphics[width=\linewidth]{\figdir edge_functions_candidate.png}
    \caption{Porte \textbf{Candidate} \textemdash{}
             arete la plus forte du reseau~:
             \texttt{v1h}$\to h^{[10]}$ ($|\varphi|_1{=}0.411$,
             deux paliers~: signature du smurfing).}
  \end{subfigure}\hfill
  \begin{subfigure}[b]{0.48\linewidth}
    \centering
    \includegraphics[width=\linewidth]{\figdir edge_functions_output.png}
    \caption{Porte \textbf{Output} \textemdash{}
             \texttt{v1h} pilote deux neurones distincts~;
             \texttt{var\_agent\_split}$\to h^{[13]}$ presente
             un pic gaussien pur tres localise.}
  \end{subfigure}

  \caption{Fonctions d'activation $\phiij(x)$ (eq.~\ref{eq:phi}) pour les
           6~aretes les plus importantes de chaque porte T-KAN ($\theta=0.01$).
           \textcolor{mkanblue}{Bleu}~: composante gaussienne FastKAN
           \cite{li2024fastkan}.
           \textcolor{lstmorange}{Orange}~: composante Fourier KAN-AD
           \cite{bozorgasl2024wav}.
           Vert~: combinaison totale $\phiij$.
           Les aretes a composante gaussienne dominante encodent des seuils
           locaux (pic ou creux)~; les aretes a composante Fourier dominante
           encodent des relations periodiques (\eg~\texttt{rho\_rupture}
           $\to h^{[15]}$ et la dynamique ATO).}
  \label{fig:edgefuncs}
\end{figure*}

La figure~\ref{fig:edgefuncs} materialise les fonctions apprises
sur les aretes survivantes. L'arete \texttt{rho\_nouveau}$\to h^{[10]}$
($|\varphi|_1=0.375$) presente une forme sigmoide douce~: plus
le ratio de nouveaux destinataires est eleve, plus la cellule
memoire est remise a zero. L'arete \texttt{v1h}$\to h^{[10]}$
($|\varphi|_1=0.411$) exhibe une non-linearite en deux paliers
caracteristique du smurfing.

\subsection{Cartes d'activation L1 (Heatmaps)}

\begin{figure*}[tbp]
  \centering
  \begin{subfigure}[b]{0.49\linewidth}
    \includegraphics[width=\linewidth]{\figdir heatmap_forget.png}
    \caption{Porte \textbf{Forget} \textemdash{} quasi-deserte~;
             \texttt{flag\_nuit}$\to h^{[9]}$ domine
             ($|\varphi|_1{\approx}0.35$).}
  \end{subfigure}\hfill
  \begin{subfigure}[b]{0.49\linewidth}
    \includegraphics[width=\linewidth]{\figdir heatmap_input.png}
    \caption{Porte \textbf{Input} \textemdash{} sparsification extreme~;
             \texttt{flag\_nuit}$\to h^{[9]}$ unique point rouge
             ($|\varphi|_1{\approx}0.22$).}
  \end{subfigure}
  \begin{subfigure}[b]{0.49\linewidth}
    \includegraphics[width=\linewidth]{\figdir heatmap_candidate.png}
    \caption{Porte \textbf{Candidate} \textemdash{} la plus dense~;
             \texttt{flag\_nuit}$\to h^{[9]}$ noir dominant
             (${\approx}0.4$), plusieurs aretes orangees
             autour de la colonne $h^{[9]}$.}
  \end{subfigure}\hfill
  \begin{subfigure}[b]{0.49\linewidth}
    \includegraphics[width=\linewidth]{\figdir heatmap_output.png}
    \caption{Porte \textbf{Output} \textemdash{} distribution diffuse
             (echelle max~$0.06$)~;
             \texttt{flag\_nuit} actif sur $h^{[9]}$ et $h^{[12]}$.}
  \end{subfigure}
  \caption{Cartes de chaleur de la norme L1 exacte $|\phiij|_1$
           pour les quatre portes T-KAN.
           Axes~: entrees ($d_h=16$ neurones d'etat + 12 variables metier)
           $\times$ sorties ($d_h=16$ neurones).
           Le signal \texttt{flag\_nuit}$\to h_{\text{out}}^{[9]}$ est le
           point le plus intense dans les quatre portes~: le neurone~9 est
           un \emph{hub} de detection des fraudes nocturnes, coherent avec
           la prevalence de ce pattern en Mobile Money CEMAC
           \cite{gsma2023}.}
  \label{fig:heatmaps}
\end{figure*}

Les heatmaps (figure~\ref{fig:heatmaps}) revelent que
\texttt{flag\_nuit}$\to h_{\text{out}}^{[9]}$ est le point le
plus intense dans les quatre portes. Le neurone $h_{\text{out}}^{[10]}$
est un \emph{hub} concentrant \texttt{rho\_nouveau} et \texttt{v1h}.
L'etat cache $h^{[13]}\to h_{\text{out}}^{[13]}$ ($|\varphi|_1=0.183$)
encode une memoire de long terme pour les schemas ATO persistants.

% =========================================================
\FloatBarrier
\section{Detection de derive conceptuelle}
% =========================================================

\begin{figure}[H]
  \centering
  \includegraphics[width=\linewidth]{\figdir drift_by_feature.png}
  \caption{Divergence JS$(P_{\text{ref}} \| P_{\text{cur}})$
           \cite{lin1991divergence} par variable entre l'entrainement
           (graine~1000) et le jeu de test (graine~1002).
           Seuil $\tau_{\text{drift}}=0.05$ (trait rouge).
           Seule \texttt{var\_agent\_split} franchit le seuil
           ($\text{JS}{=}0.051$, barre rouge)~; onze autres
           variables restent dix fois en dessous.}
  \label{fig:drift}
\end{figure}

\begin{table}[H]
\centering
\caption{Divergences Jensen-Shannon par variable
         (test vs entrainement, $\tau=0.05$).}
\label{tab:drift}
\renewcommand{\arraystretch}{1.15}
\small
\begin{tabular}{llc}
\toprule
\textbf{Variable} & \textbf{Description} & \textbf{JS} \\
\midrule
\rowcolor{alertred!20}
\texttt{var\_agent\_split} & Fractionnement agent & \textbf{0.051} \\
\texttt{v1h}               & Volume horaire       & 0.030 \\
\texttt{rho\_refund}       & Ratio remboursement  & 0.023 \\
\texttt{rho\_nouveau}      & Nouveaux destinataires & 0.008 \\
Autres (8 variables)       &                   & $< 0.001$ \\
\bottomrule
\end{tabular}
\end{table}

Seule \texttt{var\_agent\_split} depasse $\tau=0.05$
($\text{JS}=0.051$). Ce changement declenche le mecanisme de
\emph{Grid Extension} \cite{liu2024kan}, qui insere de nouveaux
centres gaussiens sans reentrainement complet.

% =========================================================
\section{Rapport d'audit symbolique COBAC}
% =========================================================

% Page paysage pour lisibilite du tableau symbolique
\clearpage
\begin{landscape}
\begin{figure}[H]
  \centering
  \includegraphics[width=0.98\linewidth,
                   height=0.62\paperwidth,
                   keepaspectratio]{\figdir symbolic_report.png}
  \caption{Rapport d'audit symbolique~: aretes survivantes par porte T-KAN,
           formule extraite par regression ($c\cdot f(ax+b)+d$),
           importance $|\phiij|_1$ et coefficient $R^2$.
           Les aretes marquees \og non-symbolifiable~(numerique)\fg{}
           ont $R^2 < 0.80$~: comportements hybrides Gaussiens+Fourier sans
           forme fermee simple, signalees explicitement conformement a la
           politique de transparence COBAC \cite{cobac2018}
           (section 2.5.5 du memoire).}
  \label{fig:symbolic}
\end{figure}
\end{landscape}
\clearpage

Pour chaque arete survivante, l'algorithme de regression symbolique
(section 2.5 du memoire) ajuste $c\cdot f(ax+b)+d$ sur une grille
de 100~points en balayant 13~familles de fonctions elementaires
($\tanh$, $\sin$, $x^2$, $e^x$, identite, etc.). Les aretes
symbolifiables ($R^2 \ge 0.80$) correspondent a des fonctions
monotones simples~; les aretes non-symbolifiables sont signalees
comme \og numeriques\fg{} conformement aux exigences COBAC
\cite{cobac2018}.

% =========================================================
\section{Explicabilite COBAC \textemdash{} analyse d'une decision}
\label{sec:decision}
% =========================================================

\begin{figure}[H]
  \centering
  \includegraphics[width=\linewidth]{\figdir decision_explanation.png}
  \caption{Explication de decision COBAC pour une transaction frauduleuse
           (score~: $\mathbf{99.97\,\%}$, verdict FRAUDE).
           Gauche~: contributions normalisees~
           \texttt{rho\_nouveau} ($0.149$) et \texttt{v1h} ($0.124$)
           dominent, suivis de \texttt{rho\_rupture} ($0.059$) et
           de la composante memorielle $h^{[13]}$ ($0.049$), attestant
           que la transaction s'inscrit dans un schema sequentiel suspect.
           Droite~: jauge de risque, seuil de decision a $50\,\%$.}
  \label{fig:decision}
\end{figure}

La figure~\ref{fig:decision} illustre le module d'explicabilite
conforme au Reglement COBAC N$^{\circ}$04/18 \cite{cobac2018}.

\paragraph{Transaction analysee.}
Score $\mathbf{99.97\,\%}$ \textemdash{} verdict FRAUDE sans ambiguite.

\paragraph{Contributions des variables.}
\[
  \Delta(\texttt{rho\_nouveau}){=}0.149,\;\;
  \Delta(\texttt{v1h}){=}0.124,\;\;
  \Delta(\texttt{rho\_rupture}){=}0.059
\]
\[
  \Delta(h^{[13]}){=}0.049,\;\;
  \Delta(\texttt{delta\_B\_orig}){=}0.041
\]
La contribution memorielle $h^{[13]}=0.049$ indique que cette
transaction ne serait pas aussi risquee isolement.

\paragraph{Valeur reglementaire.}
Ces contributions constituent la preuve algorithmique opposable
exigee par le COBAC~: l'analyste de conformite verifie que la
decision repose sur des variables metier comprehensibles, et non
sur une boite noire opaque.

% =========================================================
\section{Analyse synthetique des resultats}
\label{sec:synthese}
% =========================================================

\subsection{Deux hubs neuronaux pour deux familles de fraude}

Le neurone $h_{\text{out}}^{[9]}$ est le point le plus intense
dans les quatre heatmaps~: sa connexion dominante est
\texttt{flag\_nuit}$\to h^{[9]}$, normes L1 de $0.22$ a $0.40$
selon la porte. $h^{[9]}$ est le \emph{detecteur de fraude nocturne}
de l'architecture \cite{gsma2023}. Le neurone $h_{\text{out}}^{[10]}$
concentre \texttt{rho\_nouveau} et \texttt{v1h}~: il encode la
dynamique des comptes mule et du smurfing.

\subsection{Division fonctionnelle des portes du T-KAN}

\begin{itemize}
  \item \textbf{Forget}~: reinitialise la memoire sur \texttt{rho\_nouveau}
        eleve (comptes mule)~; 95.5\,\% elaguee.
  \item \textbf{Input}~: n'ecrit en memoire que sur
        \texttt{flag\_nuit}$\to h^{[9]}$~;
        la plus sparse (96.4\,\%).
  \item \textbf{Candidate}~: genere le contenu a ecrire~;
        la plus dense (89.1\,\%), integre la diversite des
        patterns (\texttt{v1h}, \texttt{rho\_rupture}, recurrences).
  \item \textbf{Output}~: expose l'etat cache a la projection~;
        distribution diffuse, normes $\leq 0.06$.
\end{itemize}

\subsection{Evenements de Grid Extension et stabilite}

Les perturbations synchrones aux epoques~20 et~29 sont causees par
les insertions de centres gaussiens dans la region de derive de
\texttt{var\_agent\_split} (JS$=0.051$). Le mecanisme absorbe cette
derive en 2--3~epoques sans reentrainement complet. L'absence de
derive sur les 11~autres variables confirme la precision du declenchement.

\subsection{Compromis performance / interpretabilite}

\MKAN{} se classe troisieme sur les metriques pures
($\Delta\text{MCC}\approx0.05$ vs XGBoost \cite{chen2016xgboost}),
mais est le seul modele a produire~: (i)~114~aretes avec formule
symbolique ou mention explicite de non-symbolifiabilite~;
(ii)~une decomposition en contributions normalisees par transaction~;
(iii)~un rapport d'audit opposable COBAC \cite{cobac2018}.

% =========================================================
\section{Conclusion et perspectives}
% =========================================================

\paragraph{Performance discriminative.}
\MKAN{} atteint MCC$\,{=}\,0.9268$, AUC-ROC$\,{=}\,0.9806$ sur
452\,349~fenetres de test, surpassant de $+12.8$~points la
reference externe \cite{azamuke2025} (MCC$\,{=}\,0.82$).

\paragraph{Interpretabilite structurelle.}
93.6\,\% des connexions eliminees~: 114~aretes documentent la
logique de decision. Deux hubs~: $h^{[9]}$ (fraudes nocturnes),
$h^{[10]}$ (comptes mule / smurfing).

\paragraph{Robustesse et derive.}
Un seul indicateur declenche (\texttt{var\_agent\_split},
JS$\,{=}\,0.051$), absorbe en 2--3~epoques sans reentrainement.

\paragraph{Conformite COBAC.}
Module d'explication par contributions additives~:
rapport auditable conforme au Reglement N$^{\circ}$04/18/CEMAC/UMAC/COBAC
\cite{cobac2018}, justification numerique opposable.

\paragraph{Perspectives \textemdash{} deploiement d'un dashboard temps reel.}
La prochaine etape operationnelle est un systeme de scoring temps
reel compose de quatre modules~:

\begin{enumerate}
  \item \textbf{Tampon Redis (RedisWindowBuffer).}
        Maintient les $W=10$ dernieres transactions de chaque compte
        (TTL 24h, cle \texttt{window:\{account\_id\}}). Des que la
        fenetre est pleine, le scoring est declenche automatiquement.

  \item \textbf{Inference CPU (MKANInference).}
        Chargement unique du checkpoint au demarrage, execution sur
        CPU ($<10$~ms par fenetre)~: aucune dependance a un
        accelerateur materiel. La normalisation (log-transform puis
        z-score) est appliquee dans l'ordre exact de l'entrainement.

  \item \textbf{Dashboard Dash+Plotly.}
        Interface web temps reel affichant~: jauge de risque FRAUDE/
        NORMALE, contributions normalisees par variable,
        courbe de derive JS cumulee, historique des 50 derniers scores
        du compte. Le rapport d'audit symbolique COBAC est genere a la
        demande via un bouton dedie (operation couteuse $\approx5$--$30$~s).

  \item \textbf{Mode simulation automatique.}
        Injection de transactions aleatoires (5\,\% de fraudes
        synthetiques a features extremes) sur 10~comptes simules,
        pour valider le pipeline en l'absence de donnees reelles.
\end{enumerate}

\noindent\textit{Autres axes~:}
(i)~Validation sur donnees reelles MTN Cameroun / Orange Cote d'Ivoire.
(ii)~Extension des familles symboliques ($\tanh\circ\log$, etc.)
pour ameliorer le taux de symbolification.
(iii)~Adaptation de l'architecture \MKAN{} aux transactions SWIFT
inter-banques de la BEAC.

% Bibliographie
\begin{thebibliography}{99}

\bibitem{azamuke2025}
O.~Azamuke, J.~Nakatumba-Nabende et E.~Babirye (2025).
\textit{Real-time mobile money fraud detection in sub-Saharan Africa
using ensemble learning}.
Expert Systems with Applications.

\bibitem{bozorgasl2024wav}
A.~Bozorgasl et H.~Chen (2024).
\textit{Wav-KAN: Wavelet Kolmogorov-Arnold Networks}.
arXiv:2405.12832.

\bibitem{chen2016xgboost}
T.~Chen et C.~Guestrin (2016).
\textit{XGBoost: A Scalable Tree Boosting System}.
In \textit{Proc. 22nd ACM SIGKDD}, pp. 785--794.

\bibitem{cobac2018}
COBAC (2018).
\textit{Reglement N$^{\circ}$04/18/CEMAC/UMAC/COBAC relatif aux
etablissements de paiement dans la CEMAC}.
Journal Officiel de la CEMAC.

\bibitem{gsma2023}
GSMA (2023).
\textit{State of the Industry Report on Mobile Money 2023}.
GSMA Mobile Money Programme.

\bibitem{hochreiter1997}
S.~Hochreiter et J.~Schmidhuber (1997).
Long Short-Term Memory.
\textit{Neural Computation}, 9(8), 1735--1780.

\bibitem{kingma2014adam}
D.~P.~Kingma et J.~Ba (2015).
\textit{Adam: A Method for Stochastic Optimization}.
In \textit{Proc. ICLR 2015}.

\bibitem{li2024fastkan}
Z.~Li (2024).
\textit{FastKAN: Very Fast Kolmogorov-Arnold Networks}.
arXiv:2408.07785.

\bibitem{lin1991divergence}
J.~Lin (1991).
Divergence measures based on the Shannon entropy.
\textit{IEEE Transactions on Information Theory}, 37(1), 145--151.

\bibitem{liu2024kan}
Z.~Liu, Y.~Wang, S.~Vaidya, F.~Ruehle, J.~Halverson,
M.~Soljacic, T.~Y.~Hou et M.~Tegmark (2024).
\textit{KAN: Kolmogorov-Arnold Networks}.
arXiv:2404.19756.

\bibitem{lopez2022momtsim}
A.~Lopez-Rojas et S.~Axelsson (2022).
\textit{MOMTSIM: A Multi-Agent Mobile Money Transaction Simulator}.
In \textit{Proc. IEEE Big Data}, pp. 3440--3449.

\bibitem{lundberg2017shap}
S.~M.~Lundberg et S.-I.~Lee (2017).
A unified approach to interpreting model predictions.
In \textit{Advances in Neural Information Processing Systems~30},
pp. 4765--4774.

\end{thebibliography}

\end{document}
